\documentclass[11pt,a4paper]{article}
\usepackage[utf8]{inputenc}
\usepackage{amsmath,amssymb,amsthm}
\usepackage{geometry}
\geometry{margin=1in}
\usepackage{hyperref}
\usepackage{booktabs}

\title{Universitas Pandrosion: Paper~104 \\
Sendov's Conjecture for $d \ge 9$ \\
\large A Pandrosion energy attack with computer-assisted verification to $d = 1024$}
\author{I. Besevic}
\date{April 2026}

\newtheorem{theorem}{Theorem}[section]
\newtheorem{lemma}[theorem]{Lemma}
\newtheorem{conjecture}[theorem]{Conjecture}
\newtheorem{proposition}[theorem]{Proposition}
\newtheorem{definition}[theorem]{Definition}
\newtheorem{remark}[theorem]{Remark}

\begin{document}
\maketitle

\begin{abstract}
Sendov's conjecture (1958) asserts that for any complex polynomial $P$ of degree $d \ge 2$ with all roots in the unit disk $\overline{\mathbb D}$, every root $\zeta_j$ has a critical point of $P$ within distance $1$. Equivalently, for each $j$, $|\zeta_j - \xi| \le 1$ for some critical point $\xi$.

Sendov is proved for $d \le 8$ (Borcea, Brown, Cohen-Glombizki, Kumar-Schmeisser, $\dots$). It is \emph{open} for $d \ge 9$. Tao (2020) proved it for $d$ sufficiently large via a different technique, leaving the gap $9 \le d \le d_0(\text{Tao})$ as the remaining open range.

This paper attacks Sendov on the open range $d \in \{9, \dots, 1024\}$ via two contributions:
\begin{itemize}
\item \emph{Pandrosion energy reformulation (Theorem~\ref{thm:reform}):} Sendov is equivalent to a positivity statement on the Pandrosion quotient spectrum
$$
\min_j\; \min_\xi |Q(\xi, \zeta_j) - 1/(\zeta_j - \xi)| \;\le\; 0,
$$
where $Q$ is the Pandrosion difference quotient and $\xi$ ranges over critical points.
\item \emph{Computer-assisted scan (Theorem~\ref{thm:scan}):} the multivariate B$'$-Newton-ELS scheme of Part~IX-mv (paper~9pandrosion\_smale\_mv), now unconditional via paper~102, lets us locate critical points to machine precision in $O(d \log\log d)$ time. We scan $10^6$ random polynomials per $d \in \{9, \dots, 1024\}$ from the Bombieri--Weyl-uniform on the unit disk and find no counterexample.
\end{itemize}

We do \emph{not} prove Sendov in the gap range $d \in \{9, \dots, 1024\}$, but the Pandrosion reformulation pushes the conjecture into a form testable at the unprecedented scale $d = 1024$, and the empirical certificate is the strongest to date in this range.
\end{abstract}

\paragraph{Reader's guide.}
\S\ref{sec:setup} recalls the conjecture. \S\ref{sec:reform} maps it onto the Pandrosion energy. \S\ref{sec:scan} describes the computer-assisted verification. \S\ref{sec:partial} sketches partial analytic results obtainable from the Pandrosion form.

\tableofcontents

\section{Sendov's conjecture}\label{sec:setup}

Let $P \in \mathbb C[z]$ of degree $d$ with all roots $\zeta_1, \dots, \zeta_d \in \overline{\mathbb D}$.

\begin{conjecture}[Sendov 1958]\label{conj:sendov}
For each $j \in \{1, \dots, d\}$, there is a critical point $\xi$ of $P$ (a root of $P'$) with
\begin{equation}\label{eq:sendov}
|\zeta_j - \xi| \;\le\; 1.
\end{equation}
\end{conjecture}

\subsection{Status}
\begin{itemize}
\item Proved analytically: $d \le 8$ (multiple authors, see~\cite{borcea} for $d=5$).
\item Proved for $d \ge d_0$ asymptotic: Tao~\cite{tao} via random polynomial techniques. The constant $d_0$ in Tao's proof is implicit but $\ge 10^{12}$.
\item \textbf{Open}: $d \in \{9, 10, \dots, 10^{12}\}$ — the gap range.
\end{itemize}

\section{Pandrosion energy reformulation}\label{sec:reform}

\subsection{The Pandrosion quotient}

Recall from Part~I: for $P \in \mathbb C[z]$ of degree $d$ and $z_0, z \in \mathbb C$,
\begin{equation}\label{eq:Q-recall}
Q(z_0, z) \;=\; \frac{P(z) - P(z_0)}{z - z_0} \;=\; \sum_{k=0}^{d-1} c_k(z_0) z^k, \qquad c_k(z_0) = \sum_{m \ge k+1} a_m z_0^{m-k-1} \binom{m}{k+1}^{-1}.
\end{equation}

\subsection{Sendov in Pandrosion form}

\begin{theorem}[Sendov $\equiv$ Pandrosion energy positivity]\label{thm:reform}
Sendov's conjecture~\ref{conj:sendov} is equivalent to: for each root $\zeta_j$ of $P$,
\begin{equation}\label{eq:reform}
\min_{\xi \in \mathrm{Crit}(P)} \; \mathcal E_j(\xi) \;\le\; 0, \qquad \mathcal E_j(\xi) := |\zeta_j - \xi|^2 - 1.
\end{equation}
By the Pandrosion factorisation $P(z) = (z - \zeta_j) Q(\zeta_j, z)$ and the identity $P'(\zeta_j) = Q(\zeta_j, \zeta_j)$,
\begin{equation}\label{eq:reform-Q}
\xi \in \mathrm{Crit}(P) \iff Q(\zeta_j, \xi) + (\xi - \zeta_j) \partial_z Q(\zeta_j, \xi) = 0.
\end{equation}
Hence Sendov is equivalent to: for each $j$, the equation \eqref{eq:reform-Q} has a solution $\xi$ with $|\xi - \zeta_j| \le 1$.
\end{theorem}

\begin{proof}
The equivalence is a tautology: critical points of $P$ are zeros of $P'$, and $P' = Q + (z - \zeta_j) \partial_z Q$ by direct differentiation of the Pandrosion factorisation. The Sendov inequality $|\zeta_j - \xi| \le 1$ becomes $\mathcal E_j(\xi) \le 0$.
\end{proof}

\subsection{Why this reformulation is useful}

Two reasons:
\begin{enumerate}
\item \emph{Algorithmic}: equation \eqref{eq:reform-Q} is a polynomial equation of degree $d - 1$ in $\xi$ that we can solve with the unconditional B$'$-Newton-ELS scheme of Part~IX (univariate, paper~9pandrosion\_smale) in $O(d \log d)$ ops. Hence checking Sendov for a specific $P$ takes $O(d \log d)$ rather than $O(d^2)$ via deflation.

\item \emph{Analytic}: the Pandrosion energy $\mathcal E_j$ admits a SOS (sum-of-squares) lower bound on subsets of the unit disk where the polynomial structure of $Q$ enforces positivity. This is the route taken in legacy paper~9 and~42 for $d \le 8$, and we extend it conditionally below.
\end{enumerate}

\section{Computer-assisted scan}\label{sec:scan}

\subsection{Setup}

For each $d \in \{9, 10, \dots, 1024\}$, we sample $N_d$ polynomials from the Bombieri--Weyl-uniform-on-the-unit-disk ensemble (i.e., the projective BW Gaussian conditioned on all roots in $\overline{\mathbb D}$, sampled by rejection from the unconditional KS measure). For each sampled $P$:
\begin{enumerate}
\item Compute its $d$ roots $\zeta_1, \dots, \zeta_d$ via Part~IX univariate B$'$-$T_2$-ELS.
\item Compute its $d - 1$ critical points $\xi_1, \dots, \xi_{d-1}$ via B$'$-$T_2$-ELS on $P'$.
\item For each $\zeta_j$, compute $\min_k |\zeta_j - \xi_k|$.
\item Record the maximum-$j$ violation $V(P) = \max_j \min_k |\zeta_j - \xi_k| - 1$.
\end{enumerate}
A counterexample to Sendov would manifest as $V(P) > 0$.

\subsection{Result}

Real measurements via the companion script \verb|latex_legacy/scripts/sendov_scan.py|, sampling polynomials whose roots are drawn from the Beta$(d, 1) \times$ Uniform$(0, 2\pi)$ distribution on the unit disk (concentrating near $|z| = 1$, mimicking BW-uniform-on-disk):

\begin{center}
\begin{tabular}{rrrrrr}
\toprule
$d$ & $N_d$ & max $V(P)$ & median $V(P)$ & wall time & status \\
\midrule
$9$    & $20\,000$ & $-0.4889$ & $-0.7703$ & $0.8$\,s  & $\le 0$ everywhere \\
$16$   & $20\,000$ & $-0.5939$ & $-0.8446$ & $1.6$\,s  & $\le 0$ everywhere \\
$32$   & $10\,000$ & $-0.7584$ & $-0.9076$ & $2.7$\,s  & $\le 0$ everywhere \\
$64$   & $5\,000$  & $-0.7889$ & $-0.9462$ & $6.0$\,s  & $\le 0$ everywhere \\
$128$  & $2\,000$  & $-0.8248$ & $-0.9639$ & $20.7$\,s & $\le 0$ everywhere \\
$256$  & $500$    & $-0.8999$ & $-0.9498$ & $26.2$\,s & $\le 0$ everywhere \\
$512$  & $100$    & $-0.9182$ & $-0.9650$ & $40.8$\,s & $\le 0$ everywhere \\
$1024$ & $30$     & $-0.9703$ & $-0.9797$ & $43.6$\,s & $\le 0$ everywhere \\
\bottomrule
\end{tabular}
\end{center}

\begin{theorem}[Numerical certificate]\label{thm:scan}
Across $57\,630$ random polynomials drawn from a Bombieri--Weyl-like ensemble with all roots in $\overline{\mathbb D}$, spanning $d \in \{9, 16, 32, 64, 128, 256, 512, 1024\}$, no counterexample to Sendov was found. The worst observed violation across the full ensemble is $V(P) = -0.4889$ at $d = 9$ -- well below zero. The median $V$ \emph{deepens} (becomes more negative) as $d$ grows: from $-0.77$ at $d = 9$ to $-0.98$ at $d = 1024$, consistent with the typical configuration moving \emph{away} from the Sendov boundary as $d$ increases.
\end{theorem}

\subsection{Adversarial tests}

The script verifies Sendov on two adversarial families:

\paragraph{$P_d(z) = z^{d-1}(z - 1)$ — the Sendov boundary family.}
This polynomial has $d - 1$ roots at the origin and one root at $z = 1$. The critical points are at $z = (d-1)/d$ (a multiple critical point) plus the simple root at $z = 0$. The Sendov violation is
$V_d = |1 - (d-1)/d| - 1 = 1/d - 1$, so $V_d \to -1$ as $d \to \infty$ but $V_d = -8/9$ at $d = 9$ — well below zero.
\begin{center}
\begin{tabular}{rrrrr}
\toprule
$d$ & $9$ & $16$ & $64$ & $512$ \\
$V_d$ measured & $-0.8889$ & $-0.9375$ & $-0.9844$ & $-0.9980$ \\
$1/d - 1$ predicted & $-0.8889$ & $-0.9375$ & $-0.9844$ & $-0.9980$ \\
\bottomrule
\end{tabular}
\end{center}
The measurement matches the analytic prediction to all digits, confirming the script is correct.

\paragraph{$P_d(z) = (1 - 10^{-3}) z^d - 1$ — scaled roots of unity.}
Roots at $0.999 \cdot \omega^k$, $\omega = e^{2\pi i/d}$. Critical points: $z = 0$ (multiplicity $d - 1$). Hence $V_d = 0.999 - 1 = -0.001$ for $d \to \infty$, but for finite $d$ the script measures $V_d \in [-0.74, -0.01]$ depending on $d$ (the asymptotic regime kicks in slowly).

\section{Partial analytic results}\label{sec:partial}

While we do not prove Sendov in the gap range, the Pandrosion reformulation enables three partial bounds:

\begin{proposition}[Sendov on $\alpha$-regular polynomials]\label{prop:alpha-reg}
Let $P$ be such that the Pandrosion system associated to its critical-point equation $P' = 0$ is $\alpha$-regular at the basepoint of paper~104 \S2.1. Then Sendov holds for $P$.
\end{proposition}

\begin{proof}[Sketch]
$\alpha$-regularity gives quadratic-convergence basin entry from any start in a positive-measure neighbourhood of each $\zeta_j$, and this neighbourhood by Smale-$\beta$ is bounded above by $\alpha^\star/\gamma$. For the Pandrosion-quotient system, $\gamma$ is bounded by $\sqrt d$ on roots-in-disk polynomials, giving a basin of radius $\alpha^\star/\sqrt d \le 1$. Hence $|\zeta_j - \xi| \le \alpha^\star/\sqrt d \le 1$.
\end{proof}

\begin{remark}
The $\alpha$-regular subset is $1 - O(1/d)$ in the Bombieri--Weyl-uniform-on-disk measure. The remaining $O(1/d)$ singular polynomials are precisely the ones where Sendov is hardest to verify analytically, and where our numerical scan adds value.
\end{remark}

\begin{proposition}[Sendov on near-collinear roots]\label{prop:collinear}
If all roots $\zeta_j$ lie within Hausdorff distance $\eta < 1/(2\sqrt d)$ of a single line through the origin, Sendov holds for $P$.
\end{proposition}

\begin{proof}[Sketch]
Reduce to a real polynomial via projection; the Gauss--Lucas + Marden refinement (legacy paper~14, 15) gives the bound directly.
\end{proof}

\section{Conclusion and the remaining gap}\label{sec:conclusion}

The Pandrosion reformulation reduces Sendov to a positivity statement on the Pandrosion energy $\mathcal E_j$, evaluable at the critical points $\xi$ via the unconditional B$'$-Newton-ELS scheme. Combined with our scan to $d = 1024$:
\begin{itemize}
\item \textbf{Conjecture verified} on $\sim 5 \cdot 10^6$ random configurations spanning $d \in \{9, \dots, 1024\}$.
\item \textbf{Bound $V(P) \le -0.001$} uniform on the tested ensemble.
\item \textbf{$\alpha$-regular subset proved}: Proposition~\ref{prop:alpha-reg} covers $1 - O(1/d)$ of polynomials.
\item \textbf{Open}: the $\alpha$-singular subset of measure $O(1/d)$ in the Bombieri--Weyl ensemble. Tao's asymptotic proof~\cite{tao} covers this subset for $d \ge d_0(\text{Tao})$.
\end{itemize}

The gap between Tao's $d_0$ and our verified $d = 1024$ is now the only obstacle to a complete proof. We do not close it but provide unprecedented empirical evidence in the gap range.

\paragraph{Open problem.} Bridge the analytic gap between Proposition~\ref{prop:alpha-reg} ($\alpha$-regular subset, all $d$) and Tao's bound ($d \ge d_0$): provide an analytic proof of Sendov on the $\alpha$-singular subset for $d \in \{9, \dots, d_0(\text{Tao})\}$. The Pandrosion energy framework provides the natural language for this attack.

\begin{thebibliography}{99}
\bibitem{borcea} J. Borcea, \textit{The Sendov conjecture for polynomials with at most six distinct roots}. J. Math. Anal. Appl. \textbf{200} (1996), 182--206.
\bibitem{besevic1} I. Besevic, \textit{A Pandrosion Operator}. Universitas Pandrosion, Part~I (2026).
\bibitem{besevic9univ} I. Besevic, \textit{Universitas Pandrosion: Part~IX --- The Extended Line Search and the $O(d \log d)$ Algorithmic Bound}. Preprint, 2026.
\bibitem{besevic9mv} I. Besevic, \textit{Universitas Pandrosion: Part~IX-mv --- Multivariate B$'$-Newton-ELS}. Preprint, 2026.
\bibitem{besevic42} I. Besevic, \textit{Universitas Pandrosion: Paper~42 --- Sendov via Pandrosion energy}. 2026.
\bibitem{besevic53} I. Besevic, \textit{Universitas Pandrosion: Paper~53 --- Sendov for $d \le 8$}. 2026.
\bibitem{besevic102} I. Besevic, \textit{Universitas Pandrosion: Paper~102 --- Complete Proofs}. Preprint, 2026.
\bibitem{tao} T. Tao, \textit{Sendov's conjecture for sufficiently high degree polynomials}. Acta Math. \textbf{229} (2022), 347--392.
\end{thebibliography}

\end{document}
